\scanpage{17}
% letterhead
\begin{center}
INSTITUT\\
DES HAUTES ÉTUDES SCIENTIFIQUES\\
\rule{2.3em}{0.4pt}
\end{center}

\uline{Continuation de} SGA 4 $\underline{\overline{\mathrm{IV}}}$ 4.10.6 (\inserted{\unsure{généralisation} des} 3° et 4°\unsure{)} appliqué à $S(x)\subset S_{/x}$)

5°) Soit $S_0\hookrightarrow S$ une \deleted{foncteur ($S_0$ cat. \unsure{$\mathcal{U}$|$\in\mathcal{U}$}) et supposons que} sous-catégorie pleine de $S$ satisfaisant les conditions suivantes :

(i) Si $u:X\mathbin{\unsure{-|\to}}Y$ avec $u\in\underline{M}$, $Y\in\mathrm{ob}\,S_0$, alors $X\in\mathrm{Ob}\,S_0$

[ (ii) \deleted{$\forall X\in\mathrm{Ob}\,S$, existe famille couvrante $X_i\to X$, \unsure{telle} que $X_i$ \unsure{soit}} le crible de $S$ en\uline{gen}dré par les objets de $S_0$ est couvrant, i.e. $\forall X\in\mathrm{Ob}\,S$, existe famille couvrante $X_i\to X$, avec $X_i\to Z_i\in\mathrm{Ob}\,S_0$. ]

\unsure{.} Munissons $S_0$ de la top. induite. Prouver qu'une \unsure{familles|famille} de morphismes $X_i-X$ dans $S_0$ est couvrante dans $S_0$ ssi elle l'est dans $S$\unsure{,} et que $S_0\to S$ est \uline{\unsure{cocontinu|cocontinue}}.

Soit $\underline{M}_0$ l'\unsure{ens|ens.} des $u\in\underline{M}$, \deleted{\unsure{tq}} $u:\unsure{x|X}-\unsure{y|Y}$, avec $Y\in\mathrm{Ob}\,S_0$, donc $X\in\mathrm{Ob}\,S_0$. Prouver que $\underline{M}_0\unsure{\subset|\in}\mathrm{Fl}(S_0)$ satisfait les conditions a) b) c) d) de 1°).

Procédant comme dans 3°, définir des foncteurs
\[
\mathrm{Res}:\widetilde{S}\to\widetilde{S}_0,\quad
\tikz[baseline=(p.base)]{\node[inner sep=1pt] (p) {$\mathrm{Prol}:\widetilde{S}_0\mathbin{-}\widetilde{S}_{\deleted{\unsure{0|2}}}\;?$}; \draw[gray,thin,rotate around={5:(p.center)}] (p.center) ellipse (4.9em and 1.35em);}
\]
\deleted{et prouver que}\inserted{tels que} $\mathrm{Res}$ \unsure{comm.|commute} aux limites inductives et projectives, et $\mathrm{Prol}$ est exact \unsure{:|à} \unsure{g.|q.} et pl fidèle, son image essentielle
% archive sheet number: 16
